% Standalone mathematical fragment for integration into the main article.
% Requires amsmath, amssymb, and amsthm, with lemma/theorem/corollary environments.

\section{Arithmetic endpoint certificates with binary-sized periods}

Let $x,y$ be finite words over an alphabet of signed nonzero integers.  Write
$x[i:j]$ for the half-open subword and define
\[
 \mathcal O(x,y)=\{k:1\leq k\leq\min(|x|,|y|),\quad
 x[|x|-k:|x|]=y[0:k]\}.
\]
All lengths, positions, occurrence counts, and periods are binary integers.
The words themselves are given by binary straight-line programs (SLPs).
The new construction accelerates a structural class of queries already
covered by the complete polynomial compressed matcher.  It is not a new
general complexity bound for compressed matching or for unknot recognition.

\subsection{Exact occurrence summaries}

For an explicitly represented nonempty word $a$ of length $q$, let
$S_a(w)$ be the increasing sequence of all starting positions of $a$ in $w$.
For a nonempty sequence $S=(s_1,\ldots,s_t)$ retain
\[
 \Sigma(S)=(t,s_1,s_t,\delta_{\min},\delta_{\max}),\qquad
 \delta_{\min}=\min_{i<t}(s_{i+1}-s_i),\quad
 \delta_{\max}=\max_{i<t}(s_{i+1}-s_i).
\]
Both gap entries are zero when $t=1$; the empty sequence has a separate
zero-count summary.  In particular, counts are not saturated at any fixed
threshold.

\begin{lemma}[Exact SLP occurrence summary]
Given an SLP with $g$ reachable rules and an explicit $q$-letter pattern,
the summary of every occurrence set $S_a(w)$ is computable with
$O(q^2g+g\log g)$ elementary operations and $O(qg)$ stored symbols and
integer fields.  The numerical value of the occurrence count is not a loop
bound.  The nonempty occurrence set is an arithmetic progression if and only
if its count is one or its two gap entries agree.
\end{lemma}

\begin{proof}
At a concatenation $w=uv$, the occurrences fall into three disjoint ordered
parts: those wholly in $u$, those crossing its cut, and those wholly in $v$.
The last part is shifted by $|u|$.  A crossing occurrence uses $j$ letters
of the suffix of $u$ and $q-j$ letters of the prefix of $v$, for
$1\leq j\leq q-1$.  The prefixes and suffixes of length at most $q-1$ of
each grammar word suffice to check all these candidates.  Checking at most
$q-1$ explicit words of length $q$ costs $O(q^2)$ operations per rule.

To concatenate two nonempty ordered occurrence sequences, add their counts,
keep the first position of the first and the last position of the second,
and include exactly one additional gap: the first position of the second
minus the last position of the first.  The new least and greatest gaps are
the minima and maxima of the old nonempty gap collections and this new gap.
Empty inputs and singleton gap collections require no special approximation.
This recurrence computes exact summaries by induction on the grammar.
Topological sorting contributes $O(g\log g)$ operations.  Finally all
consecutive gaps are equal exactly when their least and greatest values
agree, which is the arithmetic-progression condition.
\end{proof}

\subsection{Period-certified clipping}

An integer $d$ with $1\leq d<|P|$ is a period of $P$ if
$P[0:|P|-d]=P[d:|P|]$.  This condition can be checked by one deterministic
compressed equality call.  Neither $d$ nor the number of repetitions is
expanded, and no claim that $d$ is the least period is needed.

\begin{theorem}[Suffix-anchor clipping]
Fix $1\leq q\leq\min(|x|,|y|)$ and let $a$ be the suffix of $x$ of length
$q$.  Suppose that the complete candidate set
\[
 K=\{s+q:s\in S_a(y),\ s+q\leq |x|\}
   =\{b,b+d,\ldots,b+(r-1)d\}
\]
is a nonempty arithmetic progression with $r>1$.
Set $M=\max K$ and $P=y[0:M]$.  If $P$ has period $d$, then
\[
 \mathcal O(x,y)\cap[q,\infty)
  =K\cap\bigl[q,\operatorname{lcsuffix}(x,P)\bigr].
\]
Thus all long overlaps are obtained by one compressed longest-common-suffix
query and arithmetic clipping of a single progression.
\end{theorem}

\begin{proof}
Every overlap of length at least $q$ must align the suffix anchor of $x$
with an occurrence ending at the overlap endpoint in $y$.  It therefore
belongs to $K$.  For any $k\in K$, write $M-k=jd$ with $j\geq0$.
Repeated application of the certified period gives
$P[i]=P[i+jd]$ for $0\leq i<k$, since all intermediate indices lie in $P$.
Consequently
\[
 y[0:k]=P[0:k]=P[M-k:M].
\]
The overlap equality for this candidate is therefore exactly equality of
the length-$k$ suffixes of $x$ and $P$.  These suffixes agree if and only if
$k\leq\operatorname{lcsuffix}(x,P)$.
\end{proof}

\begin{theorem}[Prefix-anchor clipping]
Let $a=y[0:q]$ and suppose that
\[
 K=\{|x|-s:s\in S_a(x),\ |x|-s\leq |y|\}
\]
is a nonempty arithmetic progression of step $d$ and count greater than
one.  Set $M=\max K$ and $P=x[|x|-M:|x|]$.  If $P$ has period $d$, then
\[
 \mathcal O(x,y)\cap[q,\infty)
   =K\cap\bigl[q,\operatorname{lcp}(P,y)\bigr].
\]
\end{theorem}

\begin{proof}
Every sufficiently long overlap starts at an occurrence of the prefix
anchor in $x$, so it lies in $K$.  As before, $M-k$ is a nonnegative
multiple of $d$ for each candidate, and periodicity yields
$x[|x|-k:|x|]=P[M-k:M]=P[0:k]$.  Equality with $y[0:k]$ is therefore
equivalent to $k\leq\operatorname{lcp}(P,y)$.
\end{proof}

\begin{corollary}[Search by candidate rank]
Under either theorem's hypotheses, the exact predicate
\[
 f(j)=[\,x[|x|-(b+jd):|x|]=y[0:b+jd]\,],\qquad 0\leq j<r,
\]
is true on an initial segment of the candidate ranks.  Its cutoff can be
found without computing the full longest-common-prefix or suffix length.
If $T$ candidates pass and $b_{\rm bad}=r-T$, maximum/minimum tests followed
by backward doubling and bisection use
\[
 O\bigl(1+\log(1+b_{\rm bad})\bigr)
 \quad\hbox{exact substring equality calls},
\]
and hence $O(\log(r+1))$ in the worst case.  The all-true case uses one
call and the all-false case at most two.
\end{corollary}

\begin{proof}
The clipping theorems make $f$ a monotone true-prefix predicate.
First test rank $r-1$; if it passes, every rank passes.  Otherwise test rank
zero; if it fails, every rank fails.  The remaining case has known bounds
$f(0)=\mathrm{true}$ and $f(r-1)=\mathrm{false}$.
Probe backwards from $r-1$ at distances $1,2,4,\ldots$, stopping on a true
rank or on the previously checked rank zero.  A probe at distance $\delta$
passes exactly when $\delta\geq b_{\rm bad}$.  The first passing power of
two is less than $2b_{\rm bad}$, unless the search was clamped to the known
minimum.  The resulting bracket has width $O(b_{\rm bad})$.  Ordinary exact
bisection finds the last passing rank in another $O(\log(1+b_{\rm bad}))$
calls.  The case $r=2$ has no interior ranks to inspect.
\end{proof}

The final implementation uses this rank search.  The earlier full-LCP
implementation and its raw measurements are preserved as a separate ablation.
The refined comparison count is independent of the numerical period, but
each comparison and each coordinate arithmetic operation still incurs its
ordinary grammar and bit costs.

\paragraph{Empty sets, singletons, and short overlaps.}
An empty candidate set excludes every overlap of length at least $q$.
A singleton is checked by one exact substring equality, without a period
test.  The lengths $1,\ldots,q-1$ are checked separately.  Thus a successful
anchor gives a complete answer, not merely some verified witnesses.
The first candidate $b$ need not equal $d$, and the period is required only
on the longest candidate window; material outside that window is irrelevant.

\paragraph{Cropping and signs.}
Putting $L=\min(|x|,|y|)$ and replacing $x$ by its suffix of length $L$ and
$y$ by its prefix of length $L$ leaves $\mathcal O$ unchanged.  This removes
irrelevant candidates before the summaries are computed.  Reversal combined
with letterwise sign inversion is a bijection preserving equality, so the
existing free-group inverse operation turns a longest-common-suffix query
into an exact longest-common-prefix query without changing its length.
This was used by the LCP ablation; the final rank implementation directly
checks the two candidate substrings.

\paragraph{Fallback and resource bounds.}
The implementation tries $q=1,2,4,8$ in both directions, after the earlier
uniform, short-period, and sparse-endpoint cases.  A non-AP occurrence set
or a failed period certificate returns no conclusion and leaves the complete
matcher available.  A resource exception propagates; it is never interpreted
as an empty overlap set.  All query caches receive only completed results.

Let $Q$ be the largest explicit anchor length, $B$ the bit length of the
largest source-word length, and $A$ the maximum letter-label bit length.
The elementary-operation summary bound must not be interpreted as a bit-time
bound: coordinate arithmetic uses up to $B$ bits, label comparisons use up
to $A$ bits, and identifiers use $O(\log g)$ bits.  For fixed $Q=8$, the
preprocessing contributes a linear number of grammar-summary operations
apart from topological sorting, with $O(g)$ stored bounded-length buffers
and binary fields.  Exact period certification, slicing, and candidate
equalities retain their own polynomial SLP costs.  If $H$ bounds the two
input grammar heights and $c=O(1+\log(1+b_{\rm bad}))$ is the rank comparison
count, the rank search adds at most $O(Hc)$ slice rules; including short
overlaps and the bounded endpoint/period cuts gives $O(H(q+c))$ new rules.
A useful black-box bound for one successful anchor is
\[
 O\bigl(q^2g+g\log g+H(q+c)\bigr)+O(q+c+1)E(G),
\]
in elementary-operation notation, where $G$ bounds the actual intermediate
grammar size and $E$ is the selected exact-equality cost.  Thus one may use
$G=O(g+H(q+c))$ for a fresh arena containing only these inputs and the query
artifacts.  Failed earlier anchors add their explicitly bounded summary and
equality costs.  No overall linear-time claim follows from the linear summary
recurrence, and $O(\log r)$ equality calls is not itself a bit-time bound.

\subsection{A family beyond both previous thresholds}

Let $n,m$ be positive integers and
\[
 P_n=(12)^n3,\qquad W_{n,m}=P_n^m.
\]
The occurrences of the endpoint letter $3$ end at precisely
$|P_n|,2|P_n|,\ldots,m|P_n|$.  The period certificate succeeds and
$\mathcal O(W_{n,m},W_{n,m})$ is exactly this progression.  The grammar
has $O(\log n+\log m)$ rules under binary powering, while the period is
$2n+1$ and the candidate count is $m$.  Both can exceed every fixed
implementation threshold simultaneously.

The binary-alphabet variant $((21)^n1)^m$ has dense endpoint-letter
occurrences but one cyclic occurrence per block of the suffix bigram $11$;
the $q=2$ certificate yields the same progression of multiples of $2n+1$.
For $((12)^n11)^m$ with $n\geq2$, $q=4$ uses the suffix pattern $1211$;
the complete answer consists of the singleton overlap $1$ and the positive
multiples of $2n+2$ through $m(2n+2)$.

The rank refinement also handles the nonperiodic target
\[
 x=4^{2n+1}P_n^{m-1},\qquad y=P_n^m.
\]
For $m\geq3$ the exact candidate ranks have one rejected largest member.
Maximum, minimum, and immediate-predecessor tests therefore settle the whole
query in three rank predicates, regardless of the binary sizes of $n$ and
$m$.  The original full-LCP clipping may still bisect a word-length range
whose binary size includes both $\log n$ and $\log m$.

\paragraph{Scope of the improvement.}
This is a polynomial promise-case accelerator and an exact compressed
output representation inside the relator engine.  It supplies neither a
bound on the number of relator moves needed for arbitrary unknot diagrams,
nor a completeness theorem for the geometric hierarchy search, nor a
quasipolynomial bound for the complete recognition pipeline.  In particular,
the existing equality/LCP backend can still dominate phased long-period
queries even after an endpoint period certificate has succeeded.
